\section{Discussion}
\label{section:57-CD-Discussion}
We make the following observations based on the results described in section~\ref{section:56-CD-Results} and the context of the hypotheses stated in section~\ref{section:54-CD-Hypotheses}.

\textbf{The NetworkX GPU is much faster across all data shapes.}
Execution time increases proportional to the size of the graph, approximately quadratically to the number of vertices across all implementations. 
The OSM-WA dataset is the most extensive mesh graph, yet the OSM-MD mesh graph takes the longest to complete.
Both have comparable edge densities, and the somewhat irregular shape of the state of Maryland likely yields a larger diameter, emphasizing how the shape of the input graph has as much influence as the size. 
Our datasets are weighted towards mesh graphs where the number of $|E|\rightarrow|V|$.
While mesh graphs do not represent typical real-world community structures, it shows that the assertion that the Louvain algorithm and its derivatives are `near linear in practice~\cite{Blondel2008, Traag2019} is misleading on \ac{CPU} based systems.

\textbf{Execution time dominates, but ETL time is non-trivial for small graphs and \ac{GPU}s.}
Benchmarks typically focus on execution time because the execution phase is frequently the processing bottleneck, and the \ac{ETL} costs in static graph processing are assumed to be amortized over the life of the graph.
The results in figure~\ref{fig:proportions} show that for small graphs on \ac{CPU} based systems, \ac{ETL} can take up to $20\%$ of total experiment time, and for \ac{GPU} based systems can approach $75\%$ on larger graphs.
While the \ac{GPU} results in part emphasize the processing phase's speed, it highlights the overheads associated with using external accelerators that require data in a specific format for processing.
As we implement support for dynamic graphs involving many small \ac{ETL} activities over time, we expect \ac{ETL} to take more significant proportions of the total experiment runtime. 

\textbf{The Louvain Partitions implementation does not yield incremental solutions.}
Reviewing the batch latency (figure~\ref{sub@figure:batch-latency}) shows that the NetworkX implementation of Louvain Partitions does not yield an answer after each partition.
Instead, it completes the algorithm and returns a generator containing all the partitions.
To effectively measure the `benefit' of an incremental solution through batch latency, an incremental algorithm must return results in real-time, or at least in batches smaller than $99\%$ of the total runtime.

\textbf{The Louvain Partition implementation behaves inconsistently between \ac{SUT}.}
One assumption about the benefit of an incremental or approximate approach to balancing effectiveness and efficiency is that it should improve monotonically. 
That is, the longer it runs, the better the solution gets. 
Figure~\ref{fig:monotonics} shows erratic variation in the size of communities and the NMI at each step for the baseline \ac{SUT}, suggesting that we cannot make predictable trade-offs of efficiency and effectiveness on the baseline system.
The \ac{GPU} \ac{SUT} does not support the iterative louvain partition implementation. 
However, the \ac{CPU} \ac{SUT} exhibits the expected monotonic behavior, allowing users to trade-off between effectiveness and efficiency if the implementation is adjusted to print results in-stride, as an alternative approach to work within constraints if it is impractical to use a \ac{GPU}.

\textbf{Summary.}
Our small experiments show how the engineer can determine that the current baseline implementation has some issues with the effectiveness that warrant moving quickly to a different solution, that the \ac{GPU} will provide a significant speed-up on datasets at the tiny and small sizes, and that the NetworkX Louvain Partitions implementation must be updated for to trade effectiveness for efficiency.
As the \textit{NVIDIA RTXA5000} \ac{GPU} retails for approximately $USD\$2500$, above the cost of the host system without the \ac{GPU}, the engineer now has empirical solid proof to justify their system architecture for the specific problem of population segmentation. 
For all other \gai problems significant to the analysts the system will support, the engineer should repeat the benchmarking process with the respective \gai benchmark problems to gain an understanding of the trade-offs using a \ac{GPU} based approach that may require another \gai tasks. 


